\documentclass[12pt]{article}

\input{../../_assets/preambulo.tex}


\begin{document}

    % 1. Foto de fondo
    % 2. Título
    % 3. Encabezado Izquierdo
    % 4. Color de fondo
    % 5. Coord x del titulo
    % 6. Coord y del titulo
    % 7. Fecha

    
    \input{../../_assets/portada}
    \portadaExamen{ffccA4.jpg}{Curvas y\\Superficies\\Examen I}{Curvas y Superficies. Examen I}{MidnightBlue}{-8}{28}{2026}{José Juan Urrutia Milán}

    \begin{description}
        \item[Asignatura] Curvas y Superficies.
        \item[Curso Académico] 2019/20.
        \item[Grado] Doble Grado en Ingeniería Informática y Matemáticas.
        \item[Grupo] Único.
        % \item[Profesor] ---.
        \item[Descripción] Convocatoria Ordinaria.
        \item[Fecha] 17 de junio de 2020.
        \item[Duración] 3 horas.
    \end{description}
    \newpage


    % ------------------------------------
    
    \begin{ejercicio}[3 puntos]
        Sea $\alpha:I\to\mathbb{R}^3$ una curva regular que cumple $|\alpha'(t)| = 1$ y $\alpha''(t) \neq 0$, para todo $t$ del intervalo abierto $I$. Preba que $\alpha(I)$ es un arco de circunferencia si y solo si $\alpha$ tiene curvatura constante y $\alpha(I)$ está contenida en una esfera.
    \end{ejercicio}

    \begin{ejercicio}
        Considera $S = \{(x,y,z)\in \mathbb{R}^3 : -x^2+y^2+z^2 = -1, x > 0\}$.
        \begin{enumerate}[label=\alph*)]
            \item \textbf{[2 puntos]} Prueba que $S$ es una superficie regular orientable.
            \item \textbf{[2 puntos]} Calcula, en cada punto de $S$, su curvatura de Gauss y su curvatura media para una orientanción fija.
        \end{enumerate}
    \end{ejercicio}

    \begin{ejercicio}[3 puntos]
        Sea $S$ una superficie regular, compacta y conexa. Se supone que la curvatura de Gauss de $S$ cumple $K>0$ y que existe $\lm\in \mathbb{R}$ de manera que $H+\lm K= 0$, donde $H$ es la curvatura media de $S$ según una orientación fija. Prueba que $S$ es una esfera.
    \end{ejercicio}

\end{document}
